This paper presented GaitNet, a footstep planner that generates
dynamic, acyclic quadruped gaits without representing multi-leg
coordination anywhere in its architecture. At each \corlControlTickMs\,ms
control tick, GaitNet commits at most one footstep, drawn from a small
set of terrain-filtered candidates and locally refined against its own
learned value function; because a swing outlasts several such ticks,
multi-leg concurrency emerges from high-rate re-decision rather than
from an explicit schedule. On physical hardware, GaitNet (Full)
outperforms both the Legged Perceptive Stack, a state-of-the-art
perceptive NMPC baseline, and concurrency-restricted ablations of
itself (GaitNet-Single, GaitNet-Trot), achieving
\textbf{\corlHwFullSuccess\%} versus \textbf{\corlHwLPSSuccess\%} mean
success rate across matched terrain difficulties, with the gap
widening at the highest difficulties tested. Simulation
results and a swing-schedule analysis corroborate the mechanism
directly: the same policy moves one foot at a time under easy
commands and produces overlapping, non-alternating swings under
harder ones, with no change to the decision rule that produces either
behavior.

\subsection{Limitations and Future Work}

GaitNet's concurrency is capped at
\numberstringnum{\corlMaxConcurrentSwingLegs} simultaneous swing legs
by an explicit feasibility constraint requiring at least
\numberstringnum{\corlMinContactLegs} legs in ground contact; this was
a deliberate design choice for stability margin, but its effect on
performance at the extremes of terrain difficulty is untested and
worth ablating directly. The policy is also trained on a candidate
distribution, \numberstringnum{\corlTrainCandidatesPerLeg} randomly
sampled positions per leg, that differs from the distribution it is
deployed against, \numberstringnum{\corlCandidatesPerLeg} candidates
locally refined by gradient ascent; we treat this as intentional but
have not yet characterized how much the refinement step contributes
beyond what random sampling alone would achieve. Both simulation and hardware deployments depend on an
external height scan to build the terrain validity mask, and the two
platforms currently source that scan differently (a raycast height
field in simulation, the Legged Perceptive Stack's heightmap on
hardware), which is a source of sim-to-real gap we have not yet
quantified.

Future work will extend the state representation with direct
exteroceptive input, such as point clouds or height maps processed
by convolutional layers, to reduce dependence on hand-tuned terrain
thresholding, and will expand the action parameterization to include
whole-body posture and center-of-mass placement, enabling the same
high-rate re-decision principle to extend to steep inclines and
stair climbing without hard-coded elevation heuristics.
